% GENERATED FILE. Edit canonical lessons, then run npm run generate:books.
% canonical-source-hash: fnv1a64:be8e683fea04ca5f
% canonical-lessons: FR-C491-conjoint, FR-C491-separer, FR-C491-separation, FR-C491-rompre, FR-C491-rupture

\chapter{Quand ça se termine}
\label{ch:fr-quand-ca-se-termine}
\hlchaptermodality{491}

\begin{chapteropening}
\textbf{By the end of this chapter:} \emph{I can talk about breaking up and separating.}

Complete the last lesson of chapter 491: I can talk about breaking up and separating.
\end{chapteropening}

\section[le conjoint]{le conjoint — spouse, partner}

\label{lesson:FR-C491-conjoint}



\begin{quote}
Before the new one: say the French for eldest, then the French for younger.
\end{quote}

\subsection*{You'll want to know: le conjoint}
\textbf{le conjoint} — \textquotedblleft{}spouse, partner\textquotedblright{}.

\begin{quote}
Les conjoints sont invités à la fête de l'entreprise.
\end{quote}

\begin{quote}
\emph{Spouses are invited to the company party.}
\end{quote}

A woman is \emph{la conjointe}. Formal: married or living together. In Québec, \emph{le conjoint} is very common in speech.

\begin{cousinweb}
\emph{le conjoint} is the participle of \emph{conjoindre} \textquotedblleft{}to join together\textquotedblright{}, from Latin \emph{coniungere}, so \emph{coniūnctus} \textquotedblleft{}joined\textquotedblright{}, built on \emph{iungere} \textquotedblleft{}to join\textquotedblright{} (Lewis \& Short).
\end{cousinweb}

\begin{grammarlens}[title={What you've built}]
One more word for telling, explaining and describing at B1.
\end{grammarlens}

\subsection*{Your turn}
\begin{itemize}
\raggedright
  \item \emph{Say it:} \emph{le conjoint}
  \item \emph{Say it:} \emph{le conjoint}, once more
  \item \emph{Say it:} a sentence of your own with \emph{le conjoint}
  \item \emph{Recall:} say the French for eldest, then the French for younger, then say \emph{le conjoint} again
\end{itemize}

\begin{tcolorbox}[breakable,colback=teal!4,colframe=teal!35!black,title={Before you move on}]
What does le conjoint mean? (\textquotedblleft{}Spouse, partner\textquotedblright{}.) Say it once more.
\end{tcolorbox}
\section[séparer]{séparer — to separate; (se séparer) to split up}

\label{lesson:FR-C491-separer}



\begin{quote}
Before the new one: say the French for younger, then the French for spouse.
\end{quote}

\subsection*{You'll want to know: séparer}
\textbf{séparer} — \textquotedblleft{}to separate; (se séparer) to split up\textquotedblright{}.

\begin{quote}
Après dix ans de vie commune, ils se sont séparés.
\end{quote}

\begin{quote}
\emph{After ten years together, they split up.}
\end{quote}

Reflexive for couples: \emph{se séparer}.

\begin{cousinweb}
\emph{séparer} is a learned borrowing from Latin \emph{sēparāre} \textquotedblleft{}to part, to separate\textquotedblright{} (EtymDB), built from \emph{sē-} \textquotedblleft{}apart\textquotedblright{} and a verb \emph{parāre}; which \emph{parāre} is meant is not settled (Lewis \& Short).
\end{cousinweb}

\begin{grammarlens}[title={What you've built}]
One more word for telling, explaining and describing at B1.
\end{grammarlens}

\subsection*{Your turn}
\begin{itemize}
\raggedright
  \item \emph{Say it:} \emph{séparer}
  \item \emph{Say it:} \emph{séparer}, once more
  \item \emph{Say it:} a sentence of your own with \emph{séparer}
  \item \emph{Recall:} say the French for younger, then the French for spouse, then say \emph{séparer} again
\end{itemize}

\begin{tcolorbox}[breakable,colback=teal!4,colframe=teal!35!black,title={Before you move on}]
What does séparer mean? (\textquotedblleft{}To separate; (se séparer) to split up\textquotedblright{}.) Say it once more.
\end{tcolorbox}
\section[la séparation]{la séparation — separation, parting}

\label{lesson:FR-C491-separation}



\begin{quote}
Before the new one: say the French for spouse, then the French for to separate.
\end{quote}

\subsection*{You'll want to know: la séparation}
\textbf{la séparation} — \textquotedblleft{}separation, parting\textquotedblright{}.

\begin{quote}
La séparation a été difficile pour les enfants.
\end{quote}

\begin{quote}
\emph{The separation was hard for the children.}
\end{quote}

Also a parting: \emph{une séparation à la gare}.

\begin{cousinweb}
\emph{la séparation} is a learned borrowing from Latin \emph{sēparātiō} \textquotedblleft{}a separation\textquotedblright{} (EtymDB), from \emph{sēparāre} \textquotedblleft{}to part, to separate\textquotedblright{} (Lewis \& Short).
\end{cousinweb}

\begin{grammarlens}[title={What you've built}]
One more word for telling, explaining and describing at B1.
\end{grammarlens}

\subsection*{Your turn}
\begin{itemize}
\raggedright
  \item \emph{Say it:} \emph{la séparation}
  \item \emph{Say it:} \emph{la séparation}, once more
  \item \emph{Say it:} a sentence of your own with \emph{la séparation}
  \item \emph{Recall:} say the French for spouse, then the French for to separate, then say \emph{la séparation} again
\end{itemize}

\begin{tcolorbox}[breakable,colback=teal!4,colframe=teal!35!black,title={Before you move on}]
What does la séparation mean? (\textquotedblleft{}Separation, parting\textquotedblright{}.) Say it once more.
\end{tcolorbox}
\section[rompre]{rompre — to break up; to break}

\label{lesson:FR-C491-rompre}



\begin{quote}
Before the new one: say the French for to separate, then the French for separation.
\end{quote}

\subsection*{You'll want to know: rompre}
\textbf{rompre} — \textquotedblleft{}to break up; to break\textquotedblright{}.

\begin{quote}
Il a rompu avec sa copine au mois de mai.
\end{quote}

\begin{quote}
\emph{He broke up with his girlfriend in May.}
\end{quote}

\emph{Rompre avec quelqu'un} — to break up with someone. Present \emph{il rompt}.

\begin{cousinweb}
\emph{rompre} is inherited from Latin \emph{rumpere} \textquotedblleft{}to break, to burst\textquotedblright{} (EtymDB; Lewis \& Short).
\end{cousinweb}

\begin{grammarlens}[title={What you've built}]
One more word for telling, explaining and describing at B1.
\end{grammarlens}

\subsection*{Your turn}
\begin{itemize}
\raggedright
  \item \emph{Say it:} \emph{rompre}
  \item \emph{Say it:} \emph{rompre}, once more
  \item \emph{Say it:} a sentence of your own with \emph{rompre}
  \item \emph{Recall:} say the French for to separate, then the French for separation, then say \emph{rompre} again
\end{itemize}

\begin{tcolorbox}[breakable,colback=teal!4,colframe=teal!35!black,title={Before you move on}]
What does rompre mean? (\textquotedblleft{}To break up; to break\textquotedblright{}.) Say it once more.
\end{tcolorbox}
\section[la rupture]{la rupture — break-up; breaking}

\label{lesson:FR-C491-rupture}



\begin{quote}
Before the new one: say the French for separation, then the French for to break up.
\end{quote}

\subsection*{You'll want to know: la rupture}
\textbf{la rupture} — \textquotedblleft{}break-up; breaking\textquotedblright{}.

\begin{quote}
Après la rupture, elle est partie vivre à la campagne.
\end{quote}

\begin{quote}
\emph{After the break-up, she went to live in the country.}
\end{quote}

\emph{Être en rupture de stock} — to be out of stock.

\begin{cousinweb}
\emph{la rupture} is a learned borrowing from Late Latin \emph{ruptūra} \textquotedblleft{}a break\textquotedblright{}, from \emph{rumpere} \textquotedblleft{}to break\textquotedblright{} (Lewis \& Short).
\end{cousinweb}

\begin{grammarlens}[title={What you've built}]
One more word for telling, explaining and describing at B1.
\end{grammarlens}

\subsection*{Your turn}
\begin{itemize}
\raggedright
  \item \emph{Say it:} \emph{la rupture}
  \item \emph{Say it:} \emph{la rupture}, once more
  \item \emph{Say it:} a sentence of your own with \emph{la rupture}
  \item \emph{Recall:} say the French for separation, then the French for to break up, then say \emph{la rupture} again
\end{itemize}

\begin{tcolorbox}[breakable,colback=teal!4,colframe=teal!35!black,title={Before you move on}]
What does la rupture mean? (\textquotedblleft{}Break-up; breaking\textquotedblright{}.) Say it once more.
\end{tcolorbox}
